\documentclass[11pt,a4paper]{article}
\usepackage[utf8]{inputenc}
\usepackage[T1]{fontenc}
\usepackage{geometry}
\geometry{margin=1in}
\usepackage{hyperref}
\usepackage{enumitem}
\usepackage{booktabs}
\usepackage{longtable}
\usepackage{xcolor}
\usepackage{titlesec}
\usepackage{amsmath,amssymb}

\titleformat{\section}{\Large\bfseries\color{blue!70!black}}{}{0em}{}
\titleformat{\subsection}{\large\bfseries\color{blue!50!black}}{}{0em}{}
\titleformat{\subsubsection}{\normalsize\bfseries\color{blue!30!black}}{}{0em}{}

\title{\textbf{Replication Plan}\\[0.5em]
\large Constraining Cosmological Parameters with Needlet Internal Linear Combination Maps I: Analytic Power Spectrum Formalism\\[0.3em]
\normalsize OSTI ID: 2582579}
\author{Auto-generated Replication Blueprint}
\date{\today}

\begin{document}
\maketitle
\tableofcontents
\newpage

%---------------------------------------------------------------------
\section{Paper Overview}
%---------------------------------------------------------------------
This paper develops an analytic formalism for the power spectrum of Needlet Internal Linear Combination (NILC) maps used in cosmic microwave background (CMB) analysis. NILC is a component separation method that combines multi-frequency CMB observations to extract the CMB signal while minimizing foreground contamination. The paper derives analytic expressions for the residual foreground and noise biases in the NILC power spectrum, validates them against Monte Carlo simulations, and uses the formalism for cosmological parameter estimation (constraining parameters like $\Omega_b h^2$, $\Omega_c h^2$, $H_0$, $n_s$, $A_s$, $\tau$).

%---------------------------------------------------------------------
\section{Detailed Workflow}
%---------------------------------------------------------------------

\subsection{Phase 1: Code and Data Setup}
\begin{enumerate}[label=\textbf{1.\arabic*}]
  \item \textbf{Clone NILC-PS-Model repository} from GitHub.
  \item \textbf{Install dependencies}: HEALPix/healpy, CAMB or CLASS, NaMaster, NumPy, SciPy.
  \item \textbf{Download or generate sky simulations.}
\end{enumerate}

\subsection{Phase 2: Analytic NILC Power Spectrum Model}
\begin{enumerate}[label=\textbf{2.\arabic*}]
  \item \textbf{Implement needlet window functions.}
    \begin{itemize}
      \item Define needlet bands $b_j(\ell)$ covering multipole range.
      \item Implement the cosine/Gaussian needlet profiles.
    \end{itemize}
  \item \textbf{Compute ILC weights analytically.}
    \begin{itemize}
      \item Multi-frequency covariance matrix per needlet band.
      \item Minimum-variance weights: $w_j = \frac{C_j^{-1} \mathbf{a}}{\mathbf{a}^T C_j^{-1} \mathbf{a}}$ where $\mathbf{a}$ is the frequency response vector.
    \end{itemize}
  \item \textbf{Derive analytic NILC power spectrum.}
    \begin{itemize}
      \item CMB signal contribution.
      \item Residual foreground bias term.
      \item Noise bias term.
      \item ILC bias from finite sky / needlet localization.
    \end{itemize}
\end{enumerate}

\subsection{Phase 3: Monte Carlo Validation}
\begin{enumerate}[label=\textbf{3.\arabic*}]
  \item \textbf{Generate simulated multi-frequency sky maps.}
    \begin{itemize}
      \item CMB (from CAMB/CLASS $C_\ell$'s), synchrotron, dust, free-free, SZ.
      \item Instrument noise for each frequency channel (e.g., Planck-like or SO-like).
    \end{itemize}
  \item \textbf{Apply NILC pipeline} to each simulation.
  \item \textbf{Compute empirical NILC power spectrum} and compare to analytic prediction.
  \item \textbf{Verify bias formulas} against simulation average.
\end{enumerate}

\subsection{Phase 4: Cosmological Parameter Estimation}
\begin{enumerate}[label=\textbf{4.\arabic*}]
  \item \textbf{Construct the likelihood} using the analytic NILC power spectrum model.
  \item \textbf{Run MCMC} (emcee or CosmoMC) to sample the posterior of cosmological parameters.
  \item \textbf{Reproduce parameter constraints} (contour plots, marginalized posteriors).
\end{enumerate}

%---------------------------------------------------------------------
\section{Datasets}
%---------------------------------------------------------------------

\begin{longtable}{p{4.5cm} p{3.5cm} p{6cm}}
\toprule
\textbf{Dataset / Source} & \textbf{Type} & \textbf{Location / Notes} \\
\midrule
\endfirsthead
\toprule
\textbf{Dataset / Source} & \textbf{Type} & \textbf{Location / Notes} \\
\midrule
\endhead
NILC-PS-Model code + data & Repository & Check paper for GitHub URL \\
\midrule
Planck sky maps & Public CMB data & \url{https://pla.esac.esa.int/} \\
\midrule
Planck Sky Model (PSM) & Foreground simulations & \url{https://pla.esac.esa.int/} \\
\midrule
CAMB/CLASS $C_\ell$ spectra & Generated & \url{https://camb.info/} or \url{https://lesgourg.github.io/class_public/class.html} \\
\bottomrule
\end{longtable}

%---------------------------------------------------------------------
\section{Models, Tools, and Software}
%---------------------------------------------------------------------

\begin{longtable}{p{4.5cm} p{2.5cm} p{6.5cm}}
\toprule
\textbf{Tool / Library} & \textbf{Version} & \textbf{Purpose / URL} \\
\midrule
\endfirsthead
\toprule
\textbf{Tool / Library} & \textbf{Version} & \textbf{Purpose / URL} \\
\midrule
\endhead
NILC-PS-Model & latest & Paper's analysis code \newline GitHub (see paper) \\
\midrule
HEALPix / healpy & $\geq 1.15$ & Spherical harmonic transforms, sky map operations \newline \url{https://healpy.readthedocs.io/} \\
\midrule
CAMB & latest & CMB power spectrum computation \newline \url{https://camb.info/} \\
\midrule
CLASS & latest & Alternative Boltzmann solver \newline \url{https://lesgourg.github.io/class_public/class.html} \\
\midrule
NaMaster (pymaster) & latest & Pseudo-$C_\ell$ power spectrum estimation \newline \url{https://github.com/LSSTDESC/NaMaster} \\
\midrule
emcee & $\geq 3.1$ & MCMC sampler \\
\midrule
cobaya & latest & Alternative MCMC framework for cosmology \newline \url{https://cobaya.readthedocs.io/} \\
\midrule
NumPy / SciPy & $\geq 1.21$ & Numerical operations \\
\midrule
Matplotlib & $\geq 3.4$ & Power spectrum and contour plots \\
\midrule
GetDist & latest & MCMC chain analysis and triangle plots \newline \url{https://getdist.readthedocs.io/} \\
\bottomrule
\end{longtable}

\subsection{Hardware Requirements}
\begin{itemize}
  \item Monte Carlo simulations: multi-core workstation ($\geq$ 32\,GB RAM for high-resolution HEALPix maps).
  \item MCMC: can run on a single machine but benefits from parallelization.
\end{itemize}

\end{document}
